\documentclass{article}
\usepackage[utf8]{inputenc}
\usepackage[italian]{babel}

\begin{document}

\section*{Introduzione}
La \textbf{programmazione dinamica} è una tecnica basata sulle iterazioni (\textit{e non ricorsioni}).
Similmente al \textbf{\textit{divide et impera}} anche questa tecnica lavora su dei 
\textit{sotto-problemi}, ma la \textit{programmazione dinamica} conserva ogni soluzione di questi
sotto-problemi all'interno di una tabella, così da non ripetersi, a differenza della tecnica gemella;
l'uso di questa tabella evita quindi calcoli ridondanti. \\
Tra le applicazioni più conosciute delle tecnica abbiamo: \textbf{sottovettore di somma massima},
\textbf{\textit{edit distance}}, \textbf{\textit{Bellman-Ford}} e tanti altri.

\subsection*{Edit Distance (\textit{Levenshtein distance})}
L'algoritmo misura i \textbf{passaggi minimi} per trasformare una stringa $X$ in una stringa $Y$. \\
Ogni carettere della prima stringa viene confrontato con un carattere della seconda, e a ogni 
\textit{operazione} viene associato un \textbf{costo}, quest'ultimo verrà inserito nella
posizione $[x_i, y_i]$ di una matrice ausiliaria $D$, che conterrà ogni costo di ogni operazione. \\
La parte \textbf{iterativa} di questo algoritmo si può sintetizzare in questo modo:
\indent All'iterazione $(i,j)$ vengono confrontati i caratteri $x_i = y_i$, se uguali il costo 
sarà identico al passaggio precedente $D[i,j] = D[i - 1, j - 1]$, ovvero la "cella in alto a sinistra". \\
\indent Se invece $x_i \ne y_i$ si valuta che operazione applicare: % TODO %

\end{document}